%=============================================================================
% METHODOLOGY SECTION
% For inclusion in main NeurIPS 2025 document
%=============================================================================
%
% USAGE: Replace your Pipeline section with %=============================================================================
% METHODOLOGY SECTION
% For inclusion in main NeurIPS 2025 document
%=============================================================================
%
% USAGE: Replace your Pipeline section with %=============================================================================
% METHODOLOGY SECTION
% For inclusion in main NeurIPS 2025 document
%=============================================================================
%
% USAGE: Replace your Pipeline section with %=============================================================================
% METHODOLOGY SECTION
% For inclusion in main NeurIPS 2025 document
%=============================================================================
%
% USAGE: Replace your Pipeline section with \input{methodology_section}
%
% REQUIRED PACKAGES - Add to your main.tex preamble:
%   \usepackage{tikz}
%   \usetikzlibrary{positioning, arrows.meta, shapes.geometric, fit, backgrounds}
%
%=============================================================================

\section{Methodology}
\label{sec:methodology}

We present a benchmark for evaluating spatial reasoning in vision-language models that isolates relational reasoning from perceptual confounds. By providing labeled scene visualizations and structured queries with geometrically-computed ground truth, we enable precise diagnosis of model capabilities across spatial reasoning dimensions. Figure~\ref{fig:pipeline} illustrates our two-stage pipeline.

\begin{figure}[t]
\centering
\begin{tikzpicture}[
    node distance=0.4cm,
    box/.style={rectangle, draw, rounded corners=3pt, minimum width=1.8cm, minimum height=0.6cm, align=center, font=\small},
    arrow/.style={-{Stealth[length=2mm]}, thick},
    label/.style={font=\scriptsize, color=gray}
]
% Stage 1: Spatial Graph Generation
\node[box, fill=blue!15] (img) {Scene\\Image};
\node[box, fill=blue!10, right=0.5cm of img] (det) {Grounded\\SAM};
\node[box, fill=blue!10, right=0.5cm of det] (depth) {Depth\\Anything};
\node[box, fill=blue!20, right=0.5cm of depth] (graph) {Spatial\\Graph};

% Stage 2: Query Generation & Evaluation
\node[box, fill=orange!10, below=0.8cm of graph] (query) {Query\\Generation};
\node[box, fill=orange!10, left=0.5cm of query] (gt) {Ground\\Truth};
\node[box, fill=orange!15, left=0.5cm of gt] (eval) {VLM\\Evaluation};

% Arrows
\draw[arrow] (img) -- (det);
\draw[arrow] (det) -- (depth);
\draw[arrow] (depth) -- (graph);
\draw[arrow] (graph) -- (query);
\draw[arrow] (query) -- (gt);
\draw[arrow] (gt) -- (eval);

% Stage labels
\node[label, above=0.15cm of det] {Stage 1: Spatial Graph Generation};
\node[label, below=0.15cm of gt] {Stage 2: Query \& Evaluation};

\end{tikzpicture}
\caption{Pipeline overview. Stage 1 constructs a spatial graph encoding object locations and relations. Stage 2 generates queries with geometric ground truth for exact-match evaluation.}
\label{fig:pipeline}
\end{figure}

\subsection{Spatial Graph Generation}

Given an input image, we construct a spatial graph $G = (V, E)$ where nodes $V$ represent detected objects and edges $E$ encode pairwise spatial relations. Each node stores object metadata (label, bounding box, segmentation mask, depth), and each edge stores the spatial relation between the connected objects.

\paragraph{Object detection and segmentation.} We employ Grounded-SAM, combining Grounding DINO for open-vocabulary detection with SAM for precise mask generation. Given a set of object category prompts, the model returns bounding boxes and segmentation masks for all detected instances. Each object receives a unique identifier by appending an instance index to its category label (e.g., \texttt{chair\_0}, \texttt{chair\_1}). We filter detections by confidence threshold and remove heavily overlapping boxes via non-maximum suppression.

\paragraph{Depth estimation.} We apply Depth Anything V2 to obtain a dense relative depth map $D \in \mathbb{R}^{H \times W}$, where lower values indicate closer proximity to the camera. For each detected object with segmentation mask $M_i$, we compute its representative depth as the median depth value within the mask: $d_i = \text{median}(\{D[p] : p \in M_i\})$. The median provides robustness to depth discontinuities at object boundaries.

\paragraph{Spatial relation computation.} For each ordered pair of objects $(o_i, o_j)$, we compute spatial relations along four axes by comparing their bounding box centers $(c_x, c_y)$ and depth values $d$:

\begin{itemize}[leftmargin=1.5em, itemsep=2pt, topsep=3pt]
    \item \textbf{Horizontal:} $o_i$ is \textit{left} of $o_j$ if $c_x^{(i)} < c_x^{(j)}$, else \textit{right}.
    \item \textbf{Vertical:} $o_i$ is \textit{above} $o_j$ if $c_y^{(i)} < c_y^{(j)}$, else \textit{below}.
    \item \textbf{Depth:} $o_i$ is \textit{in front of} $o_j$ if $d_i < d_j$ (closer to camera), else \textit{behind}.
    \item \textbf{Saliency:} $o_i$ is \textit{foreground} relative to $o_j$ if it is both closer (lower depth) and occupies a larger screen area, else \textit{background}.
\end{itemize}

Relations are stored as directed edges in the spatial graph, enabling efficient lookup during query generation.

\subsection{Query Taxonomy}

We organize queries along three orthogonal dimensions that isolate specific reasoning competencies.

\paragraph{Egocentric vs.\ Allocentric.} \textit{Egocentric} queries use the viewer's (camera's) reference frame: ``Is the chair to the left or right of the table?'' \textit{Allocentric} queries require perspective transformation by specifying an anchor object and facing direction: ``Standing at the chair facing the window, is the lamp to your left or right?'' This dimension tests mental rotation and reference frame transformation---a capability known to be challenging for both humans and models.

\paragraph{Viewer-centered vs.\ Object-to-object.} \textit{Viewer-centered} queries ask about object positions relative to the observer (``Which object is in the foreground?''), while \textit{object-to-object} queries ask about relations between two specific objects (``Is object A above or below object B?'').

\paragraph{Response format.} We include two formats: \textit{binary choice} (select between two relation labels, e.g., left vs.\ right) and \textit{multiple choice} (select among several labels, e.g., left/right/front/behind).

\subsection{Ground Truth Computation}

Ground truth answers are computed geometrically from the spatial graph, ensuring objective and reproducible evaluation without human annotation.

\paragraph{Egocentric relations.} Relations are computed directly in image coordinates. For a query about objects $o_i$ and $o_j$, we retrieve the precomputed relation from the spatial graph edge $(o_i, o_j)$.

\paragraph{Allocentric relations.} We transform the coordinate system to the anchor object's reference frame. Given anchor object $A$ at position $\mathbf{p}_A$ and facing direction toward object $F$ at $\mathbf{p}_F$, we define a local coordinate system with forward vector $\hat{\mathbf{f}} = (\mathbf{p}_F - \mathbf{p}_A) / \|\mathbf{p}_F - \mathbf{p}_A\|$. For a target object $T$, we compute its position relative to the anchor: $\mathbf{r} = \mathbf{p}_T - \mathbf{p}_A$. Left/right is determined by the sign of $\hat{\mathbf{f}} \times \mathbf{r}$ (cross product), and front/behind by the sign of $\hat{\mathbf{f}} \cdot \mathbf{r}$ (dot product).

\subsection{Evaluation Protocol}

\paragraph{Input format.} Models receive the scene image with overlaid bounding boxes and text labels, plus a natural language query. The visual annotation ensures models can identify referenced objects without requiring object detection capabilities, isolating spatial reasoning from perception.

\paragraph{Metrics.} We report accuracy as the fraction of correctly answered queries, with breakdowns by task type (egocentric/allocentric), frame type (viewer-centered/object-to-object), and relation axis (horizontal/vertical/depth) for fine-grained diagnosis of model capabilities.

%
% REQUIRED PACKAGES - Add to your main.tex preamble:
%   \usepackage{tikz}
%   \usetikzlibrary{positioning, arrows.meta, shapes.geometric, fit, backgrounds}
%
%=============================================================================

\section{Methodology}
\label{sec:methodology}

We present a benchmark for evaluating spatial reasoning in vision-language models that isolates relational reasoning from perceptual confounds. By providing labeled scene visualizations and structured queries with geometrically-computed ground truth, we enable precise diagnosis of model capabilities across spatial reasoning dimensions. Figure~\ref{fig:pipeline} illustrates our two-stage pipeline.

\begin{figure}[t]
\centering
\begin{tikzpicture}[
    node distance=0.4cm,
    box/.style={rectangle, draw, rounded corners=3pt, minimum width=1.8cm, minimum height=0.6cm, align=center, font=\small},
    arrow/.style={-{Stealth[length=2mm]}, thick},
    label/.style={font=\scriptsize, color=gray}
]
% Stage 1: Spatial Graph Generation
\node[box, fill=blue!15] (img) {Scene\\Image};
\node[box, fill=blue!10, right=0.5cm of img] (det) {Grounded\\SAM};
\node[box, fill=blue!10, right=0.5cm of det] (depth) {Depth\\Anything};
\node[box, fill=blue!20, right=0.5cm of depth] (graph) {Spatial\\Graph};

% Stage 2: Query Generation & Evaluation
\node[box, fill=orange!10, below=0.8cm of graph] (query) {Query\\Generation};
\node[box, fill=orange!10, left=0.5cm of query] (gt) {Ground\\Truth};
\node[box, fill=orange!15, left=0.5cm of gt] (eval) {VLM\\Evaluation};

% Arrows
\draw[arrow] (img) -- (det);
\draw[arrow] (det) -- (depth);
\draw[arrow] (depth) -- (graph);
\draw[arrow] (graph) -- (query);
\draw[arrow] (query) -- (gt);
\draw[arrow] (gt) -- (eval);

% Stage labels
\node[label, above=0.15cm of det] {Stage 1: Spatial Graph Generation};
\node[label, below=0.15cm of gt] {Stage 2: Query \& Evaluation};

\end{tikzpicture}
\caption{Pipeline overview. Stage 1 constructs a spatial graph encoding object locations and relations. Stage 2 generates queries with geometric ground truth for exact-match evaluation.}
\label{fig:pipeline}
\end{figure}

\subsection{Spatial Graph Generation}

Given an input image, we construct a spatial graph $G = (V, E)$ where nodes $V$ represent detected objects and edges $E$ encode pairwise spatial relations. Each node stores object metadata (label, bounding box, segmentation mask, depth), and each edge stores the spatial relation between the connected objects.

\paragraph{Object detection and segmentation.} We employ Grounded-SAM, combining Grounding DINO for open-vocabulary detection with SAM for precise mask generation. Given a set of object category prompts, the model returns bounding boxes and segmentation masks for all detected instances. Each object receives a unique identifier by appending an instance index to its category label (e.g., \texttt{chair\_0}, \texttt{chair\_1}). We filter detections by confidence threshold and remove heavily overlapping boxes via non-maximum suppression.

\paragraph{Depth estimation.} We apply Depth Anything V2 to obtain a dense relative depth map $D \in \mathbb{R}^{H \times W}$, where lower values indicate closer proximity to the camera. For each detected object with segmentation mask $M_i$, we compute its representative depth as the median depth value within the mask: $d_i = \text{median}(\{D[p] : p \in M_i\})$. The median provides robustness to depth discontinuities at object boundaries.

\paragraph{Spatial relation computation.} For each ordered pair of objects $(o_i, o_j)$, we compute spatial relations along four axes by comparing their bounding box centers $(c_x, c_y)$ and depth values $d$:

\begin{itemize}[leftmargin=1.5em, itemsep=2pt, topsep=3pt]
    \item \textbf{Horizontal:} $o_i$ is \textit{left} of $o_j$ if $c_x^{(i)} < c_x^{(j)}$, else \textit{right}.
    \item \textbf{Vertical:} $o_i$ is \textit{above} $o_j$ if $c_y^{(i)} < c_y^{(j)}$, else \textit{below}.
    \item \textbf{Depth:} $o_i$ is \textit{in front of} $o_j$ if $d_i < d_j$ (closer to camera), else \textit{behind}.
    \item \textbf{Saliency:} $o_i$ is \textit{foreground} relative to $o_j$ if it is both closer (lower depth) and occupies a larger screen area, else \textit{background}.
\end{itemize}

Relations are stored as directed edges in the spatial graph, enabling efficient lookup during query generation.

\subsection{Query Taxonomy}

We organize queries along three orthogonal dimensions that isolate specific reasoning competencies.

\paragraph{Egocentric vs.\ Allocentric.} \textit{Egocentric} queries use the viewer's (camera's) reference frame: ``Is the chair to the left or right of the table?'' \textit{Allocentric} queries require perspective transformation by specifying an anchor object and facing direction: ``Standing at the chair facing the window, is the lamp to your left or right?'' This dimension tests mental rotation and reference frame transformation---a capability known to be challenging for both humans and models.

\paragraph{Viewer-centered vs.\ Object-to-object.} \textit{Viewer-centered} queries ask about object positions relative to the observer (``Which object is in the foreground?''), while \textit{object-to-object} queries ask about relations between two specific objects (``Is object A above or below object B?'').

\paragraph{Response format.} We include two formats: \textit{binary choice} (select between two relation labels, e.g., left vs.\ right) and \textit{multiple choice} (select among several labels, e.g., left/right/front/behind).

\subsection{Ground Truth Computation}

Ground truth answers are computed geometrically from the spatial graph, ensuring objective and reproducible evaluation without human annotation.

\paragraph{Egocentric relations.} Relations are computed directly in image coordinates. For a query about objects $o_i$ and $o_j$, we retrieve the precomputed relation from the spatial graph edge $(o_i, o_j)$.

\paragraph{Allocentric relations.} We transform the coordinate system to the anchor object's reference frame. Given anchor object $A$ at position $\mathbf{p}_A$ and facing direction toward object $F$ at $\mathbf{p}_F$, we define a local coordinate system with forward vector $\hat{\mathbf{f}} = (\mathbf{p}_F - \mathbf{p}_A) / \|\mathbf{p}_F - \mathbf{p}_A\|$. For a target object $T$, we compute its position relative to the anchor: $\mathbf{r} = \mathbf{p}_T - \mathbf{p}_A$. Left/right is determined by the sign of $\hat{\mathbf{f}} \times \mathbf{r}$ (cross product), and front/behind by the sign of $\hat{\mathbf{f}} \cdot \mathbf{r}$ (dot product).

\subsection{Evaluation Protocol}

\paragraph{Input format.} Models receive the scene image with overlaid bounding boxes and text labels, plus a natural language query. The visual annotation ensures models can identify referenced objects without requiring object detection capabilities, isolating spatial reasoning from perception.

\paragraph{Metrics.} We report accuracy as the fraction of correctly answered queries, with breakdowns by task type (egocentric/allocentric), frame type (viewer-centered/object-to-object), and relation axis (horizontal/vertical/depth) for fine-grained diagnosis of model capabilities.

%
% REQUIRED PACKAGES - Add to your main.tex preamble:
%   \usepackage{tikz}
%   \usetikzlibrary{positioning, arrows.meta, shapes.geometric, fit, backgrounds}
%
%=============================================================================

\section{Methodology}
\label{sec:methodology}

We present a benchmark for evaluating spatial reasoning in vision-language models that isolates relational reasoning from perceptual confounds. By providing labeled scene visualizations and structured queries with geometrically-computed ground truth, we enable precise diagnosis of model capabilities across spatial reasoning dimensions. Figure~\ref{fig:pipeline} illustrates our two-stage pipeline.

\begin{figure}[t]
\centering
\begin{tikzpicture}[
    node distance=0.4cm,
    box/.style={rectangle, draw, rounded corners=3pt, minimum width=1.8cm, minimum height=0.6cm, align=center, font=\small},
    arrow/.style={-{Stealth[length=2mm]}, thick},
    label/.style={font=\scriptsize, color=gray}
]
% Stage 1: Spatial Graph Generation
\node[box, fill=blue!15] (img) {Scene\\Image};
\node[box, fill=blue!10, right=0.5cm of img] (det) {Grounded\\SAM};
\node[box, fill=blue!10, right=0.5cm of det] (depth) {Depth\\Anything};
\node[box, fill=blue!20, right=0.5cm of depth] (graph) {Spatial\\Graph};

% Stage 2: Query Generation & Evaluation
\node[box, fill=orange!10, below=0.8cm of graph] (query) {Query\\Generation};
\node[box, fill=orange!10, left=0.5cm of query] (gt) {Ground\\Truth};
\node[box, fill=orange!15, left=0.5cm of gt] (eval) {VLM\\Evaluation};

% Arrows
\draw[arrow] (img) -- (det);
\draw[arrow] (det) -- (depth);
\draw[arrow] (depth) -- (graph);
\draw[arrow] (graph) -- (query);
\draw[arrow] (query) -- (gt);
\draw[arrow] (gt) -- (eval);

% Stage labels
\node[label, above=0.15cm of det] {Stage 1: Spatial Graph Generation};
\node[label, below=0.15cm of gt] {Stage 2: Query \& Evaluation};

\end{tikzpicture}
\caption{Pipeline overview. Stage 1 constructs a spatial graph encoding object locations and relations. Stage 2 generates queries with geometric ground truth for exact-match evaluation.}
\label{fig:pipeline}
\end{figure}

\subsection{Spatial Graph Generation}

Given an input image, we construct a spatial graph $G = (V, E)$ where nodes $V$ represent detected objects and edges $E$ encode pairwise spatial relations. Each node stores object metadata (label, bounding box, segmentation mask, depth), and each edge stores the spatial relation between the connected objects.

\paragraph{Object detection and segmentation.} We employ Grounded-SAM, combining Grounding DINO for open-vocabulary detection with SAM for precise mask generation. Given a set of object category prompts, the model returns bounding boxes and segmentation masks for all detected instances. Each object receives a unique identifier by appending an instance index to its category label (e.g., \texttt{chair\_0}, \texttt{chair\_1}). We filter detections by confidence threshold and remove heavily overlapping boxes via non-maximum suppression.

\paragraph{Depth estimation.} We apply Depth Anything V2 to obtain a dense relative depth map $D \in \mathbb{R}^{H \times W}$, where lower values indicate closer proximity to the camera. For each detected object with segmentation mask $M_i$, we compute its representative depth as the median depth value within the mask: $d_i = \text{median}(\{D[p] : p \in M_i\})$. The median provides robustness to depth discontinuities at object boundaries.

\paragraph{Spatial relation computation.} For each ordered pair of objects $(o_i, o_j)$, we compute spatial relations along four axes by comparing their bounding box centers $(c_x, c_y)$ and depth values $d$:

\begin{itemize}[leftmargin=1.5em, itemsep=2pt, topsep=3pt]
    \item \textbf{Horizontal:} $o_i$ is \textit{left} of $o_j$ if $c_x^{(i)} < c_x^{(j)}$, else \textit{right}.
    \item \textbf{Vertical:} $o_i$ is \textit{above} $o_j$ if $c_y^{(i)} < c_y^{(j)}$, else \textit{below}.
    \item \textbf{Depth:} $o_i$ is \textit{in front of} $o_j$ if $d_i < d_j$ (closer to camera), else \textit{behind}.
    \item \textbf{Saliency:} $o_i$ is \textit{foreground} relative to $o_j$ if it is both closer (lower depth) and occupies a larger screen area, else \textit{background}.
\end{itemize}

Relations are stored as directed edges in the spatial graph, enabling efficient lookup during query generation.

\subsection{Query Taxonomy}

We organize queries along three orthogonal dimensions that isolate specific reasoning competencies.

\paragraph{Egocentric vs.\ Allocentric.} \textit{Egocentric} queries use the viewer's (camera's) reference frame: ``Is the chair to the left or right of the table?'' \textit{Allocentric} queries require perspective transformation by specifying an anchor object and facing direction: ``Standing at the chair facing the window, is the lamp to your left or right?'' This dimension tests mental rotation and reference frame transformation---a capability known to be challenging for both humans and models.

\paragraph{Viewer-centered vs.\ Object-to-object.} \textit{Viewer-centered} queries ask about object positions relative to the observer (``Which object is in the foreground?''), while \textit{object-to-object} queries ask about relations between two specific objects (``Is object A above or below object B?'').

\paragraph{Response format.} We include two formats: \textit{binary choice} (select between two relation labels, e.g., left vs.\ right) and \textit{multiple choice} (select among several labels, e.g., left/right/front/behind).

\subsection{Ground Truth Computation}

Ground truth answers are computed geometrically from the spatial graph, ensuring objective and reproducible evaluation without human annotation.

\paragraph{Egocentric relations.} Relations are computed directly in image coordinates. For a query about objects $o_i$ and $o_j$, we retrieve the precomputed relation from the spatial graph edge $(o_i, o_j)$.

\paragraph{Allocentric relations.} We transform the coordinate system to the anchor object's reference frame. Given anchor object $A$ at position $\mathbf{p}_A$ and facing direction toward object $F$ at $\mathbf{p}_F$, we define a local coordinate system with forward vector $\hat{\mathbf{f}} = (\mathbf{p}_F - \mathbf{p}_A) / \|\mathbf{p}_F - \mathbf{p}_A\|$. For a target object $T$, we compute its position relative to the anchor: $\mathbf{r} = \mathbf{p}_T - \mathbf{p}_A$. Left/right is determined by the sign of $\hat{\mathbf{f}} \times \mathbf{r}$ (cross product), and front/behind by the sign of $\hat{\mathbf{f}} \cdot \mathbf{r}$ (dot product).

\subsection{Evaluation Protocol}

\paragraph{Input format.} Models receive the scene image with overlaid bounding boxes and text labels, plus a natural language query. The visual annotation ensures models can identify referenced objects without requiring object detection capabilities, isolating spatial reasoning from perception.

\paragraph{Metrics.} We report accuracy as the fraction of correctly answered queries, with breakdowns by task type (egocentric/allocentric), frame type (viewer-centered/object-to-object), and relation axis (horizontal/vertical/depth) for fine-grained diagnosis of model capabilities.

%
% REQUIRED PACKAGES - Add to your main.tex preamble:
%   \usepackage{tikz}
%   \usetikzlibrary{positioning, arrows.meta, shapes.geometric, fit, backgrounds}
%
%=============================================================================

\section{Methodology}
\label{sec:methodology}

We present a benchmark for evaluating spatial reasoning in vision-language models that isolates relational reasoning from perceptual confounds. By providing labeled scene visualizations and structured queries with geometrically-computed ground truth, we enable precise diagnosis of model capabilities across spatial reasoning dimensions. Figure~\ref{fig:pipeline} illustrates our two-stage pipeline.

\begin{figure}[t]
\centering
\begin{tikzpicture}[
    node distance=0.4cm,
    box/.style={rectangle, draw, rounded corners=3pt, minimum width=1.8cm, minimum height=0.6cm, align=center, font=\small},
    arrow/.style={-{Stealth[length=2mm]}, thick},
    label/.style={font=\scriptsize, color=gray}
]
% Stage 1: Spatial Graph Generation
\node[box, fill=blue!15] (img) {Scene\\Image};
\node[box, fill=blue!10, right=0.5cm of img] (det) {Grounded\\SAM};
\node[box, fill=blue!10, right=0.5cm of det] (depth) {Depth\\Anything};
\node[box, fill=blue!20, right=0.5cm of depth] (graph) {Spatial\\Graph};

% Stage 2: Query Generation & Evaluation
\node[box, fill=orange!10, below=0.8cm of graph] (query) {Query\\Generation};
\node[box, fill=orange!10, left=0.5cm of query] (gt) {Ground\\Truth};
\node[box, fill=orange!15, left=0.5cm of gt] (eval) {VLM\\Evaluation};

% Arrows
\draw[arrow] (img) -- (det);
\draw[arrow] (det) -- (depth);
\draw[arrow] (depth) -- (graph);
\draw[arrow] (graph) -- (query);
\draw[arrow] (query) -- (gt);
\draw[arrow] (gt) -- (eval);

% Stage labels
\node[label, above=0.15cm of det] {Stage 1: Spatial Graph Generation};
\node[label, below=0.15cm of gt] {Stage 2: Query \& Evaluation};

\end{tikzpicture}
\caption{Pipeline overview. Stage 1 constructs a spatial graph encoding object locations and relations. Stage 2 generates queries with geometric ground truth for exact-match evaluation.}
\label{fig:pipeline}
\end{figure}

\subsection{Spatial Graph Generation}

Given an input image, we construct a spatial graph $G = (V, E)$ where nodes $V$ represent detected objects and edges $E$ encode pairwise spatial relations. Each node stores object metadata (label, bounding box, segmentation mask, depth), and each edge stores the spatial relation between the connected objects.

\paragraph{Object detection and segmentation.} We employ Grounded-SAM, combining Grounding DINO for open-vocabulary detection with SAM for precise mask generation. Given a set of object category prompts, the model returns bounding boxes and segmentation masks for all detected instances. Each object receives a unique identifier by appending an instance index to its category label (e.g., \texttt{chair\_0}, \texttt{chair\_1}). We filter detections by confidence threshold and remove heavily overlapping boxes via non-maximum suppression.

\paragraph{Depth estimation.} We apply Depth Anything V2 to obtain a dense relative depth map $D \in \mathbb{R}^{H \times W}$, where lower values indicate closer proximity to the camera. For each detected object with segmentation mask $M_i$, we compute its representative depth as the median depth value within the mask: $d_i = \text{median}(\{D[p] : p \in M_i\})$. The median provides robustness to depth discontinuities at object boundaries.

\paragraph{Spatial relation computation.} For each ordered pair of objects $(o_i, o_j)$, we compute spatial relations along four axes by comparing their bounding box centers $(c_x, c_y)$ and depth values $d$:

\begin{itemize}[leftmargin=1.5em, itemsep=2pt, topsep=3pt]
    \item \textbf{Horizontal:} $o_i$ is \textit{left} of $o_j$ if $c_x^{(i)} < c_x^{(j)}$, else \textit{right}.
    \item \textbf{Vertical:} $o_i$ is \textit{above} $o_j$ if $c_y^{(i)} < c_y^{(j)}$, else \textit{below}.
    \item \textbf{Depth:} $o_i$ is \textit{in front of} $o_j$ if $d_i < d_j$ (closer to camera), else \textit{behind}.
    \item \textbf{Saliency:} $o_i$ is \textit{foreground} relative to $o_j$ if it is both closer (lower depth) and occupies a larger screen area, else \textit{background}.
\end{itemize}

Relations are stored as directed edges in the spatial graph, enabling efficient lookup during query generation.

\subsection{Query Taxonomy}

We organize queries along three orthogonal dimensions that isolate specific reasoning competencies.

\paragraph{Egocentric vs.\ Allocentric.} \textit{Egocentric} queries use the viewer's (camera's) reference frame: ``Is the chair to the left or right of the table?'' \textit{Allocentric} queries require perspective transformation by specifying an anchor object and facing direction: ``Standing at the chair facing the window, is the lamp to your left or right?'' This dimension tests mental rotation and reference frame transformation---a capability known to be challenging for both humans and models.

\paragraph{Viewer-centered vs.\ Object-to-object.} \textit{Viewer-centered} queries ask about object positions relative to the observer (``Which object is in the foreground?''), while \textit{object-to-object} queries ask about relations between two specific objects (``Is object A above or below object B?'').

\paragraph{Response format.} We include two formats: \textit{binary choice} (select between two relation labels, e.g., left vs.\ right) and \textit{multiple choice} (select among several labels, e.g., left/right/front/behind).

\subsection{Ground Truth Computation}

Ground truth answers are computed geometrically from the spatial graph, ensuring objective and reproducible evaluation without human annotation.

\paragraph{Egocentric relations.} Relations are computed directly in image coordinates. For a query about objects $o_i$ and $o_j$, we retrieve the precomputed relation from the spatial graph edge $(o_i, o_j)$.

\paragraph{Allocentric relations.} We transform the coordinate system to the anchor object's reference frame. Given anchor object $A$ at position $\mathbf{p}_A$ and facing direction toward object $F$ at $\mathbf{p}_F$, we define a local coordinate system with forward vector $\hat{\mathbf{f}} = (\mathbf{p}_F - \mathbf{p}_A) / \|\mathbf{p}_F - \mathbf{p}_A\|$. For a target object $T$, we compute its position relative to the anchor: $\mathbf{r} = \mathbf{p}_T - \mathbf{p}_A$. Left/right is determined by the sign of $\hat{\mathbf{f}} \times \mathbf{r}$ (cross product), and front/behind by the sign of $\hat{\mathbf{f}} \cdot \mathbf{r}$ (dot product).

\subsection{Evaluation Protocol}

\paragraph{Input format.} Models receive the scene image with overlaid bounding boxes and text labels, plus a natural language query. The visual annotation ensures models can identify referenced objects without requiring object detection capabilities, isolating spatial reasoning from perception.

\paragraph{Metrics.} We report accuracy as the fraction of correctly answered queries, with breakdowns by task type (egocentric/allocentric), frame type (viewer-centered/object-to-object), and relation axis (horizontal/vertical/depth) for fine-grained diagnosis of model capabilities.
